The canonical image consists of the primes $\mathfrak q\subseteq\mathfrak p$. Such a prime lies in every basic neighborhood $D(f)$ of $\mathfrak p$, since $f\notin\mathfrak p$ implies $f\notin\mathfrak q$. Conversely, if $\mathfrak q\nsubseteq\mathfrak p$, choose $f\in\mathfrak q\setminus\mathfrak p$. Then $D(f)$ contains $\mathfrak p$ but not $\mathfrak q$. Since the basic opens form a basis, this proves the assertion.
